\section{Quotient pathologies for Picard schemes}
\subsection{Quotients: the scope of the product formula}
\paragraph{D0 (Refutation of R6 as stated; correction of scope).}
At 57:286 the transcription explicitly assumes
``supposons que $A_{\Gamma}=0$ \dots et $\Gamma$ opère trivial[ement]
sur $\operatorname{Pic}^{\tau}_{B/S}$'' before writing the full Picard product.
Those hypotheses were omitted from R6 in the research prompt.  Without them
the full product formula is false: take $\Gamma=1$ and $A=B=E$, an elliptic
curve.  The Poincaré line bundle on $E\times E$ is not an exterior product:
its restriction to $E\times\{b\}$ varies with $b$, whereas that of an
exterior product has constant class.  Thus the missing correspondence
factor is essential.  This is a correction of the prompt, not of the
transcribed manuscript.

We use the following form, sufficient for all examples below.  Let $A,B$
be abelian schemes over a locally noetherian base $S$, and let a constant
finite group $\Gamma$ act on $A$ by group automorphisms and on $B$ by
translations through a subgroup $\Gamma_S\hookrightarrow B$.  Set
$Y=B/\Gamma$ and $X=(A\times_S B)/\Gamma$, with $f:X\to Y$.  Then, as fppf
group sheaves, and hence as their representing group algebraic spaces,
\begin{equation}\label{dgj:pic-tau}
 \operatorname{Pic}^{\tau}_{X/S}
 \simeq Y^\vee\times_S(A^\vee)^\Gamma.
\end{equation}
Here $\operatorname{Pic}^{\tau}$ means the union of geometric components
whose classes in the Néron--Severi group are torsion.  The action on the
dual is pullback; replacing all actions by inverse pullback gives the
same fixed subgroup.

For clarity we give the descent argument.  The zero of $A$ induces a
section $j:Y\to X$, so $j^*f^*=1$ splits off $\operatorname{Pic}_{Y/S}$.
An algebraically trivial line bundle on an abelian scheme, rigidified at
zero, has no nontrivial rigidified automorphism.  Consequently a fixed
point of $A^\vee$ has a unique rigidified $\Gamma$-linearization and its
pullback to $A\times B$ descends to $X$, trivially on $j(Y)$.  This
constructs a map to the relative Picard sheaf.  It lands in
$\operatorname{Pic}^{\tau}$: on each geometric fibre, its pullback
under the finite étale quotient $q:A\times B\to X$ is algebraically
trivial, while the norm identity
$\operatorname{Nm}_q(q^*N)=N^{\otimes|\Gamma|}$ shows that its
Néron--Severi class is torsion.  Norm preserves algebraic triviality,
since it carries connected families of line bundles to connected
families and sends the trivial bundle to the trivial bundle.
Conversely, a $\tau$-class pulls back to
an algebraically trivial class on $A\times B$: the Néron--Severi group of
an abelian variety is torsion-free.  The product duality
$(A\times B)^\vee=A^\vee\times B^\vee$ then shows that, after subtracting
the pullback of its restriction to $j(Y)$, it is a constant class from
$A^\vee$.  Descent makes this class fixed.  A discrepancy in descent is
a character of $\Gamma$.  Its restriction along $j$ is the line
bundle descended along $B\to Y$ by that character.  The resulting
map $\Gamma^D\hookrightarrow Y^\vee$ is the injective kernel map in
the dual isogeny sequence for $B\to Y$; thus restriction along $j$
removes the discrepancy even on nonreduced test schemes.  These
constructions commute with base change and
are inverse.  The same argument applies fppf locally to classes rather
than line bundles, so no unmentioned universal line bundle is required.

The quotient map $A\times B\to X$ is finite étale and surjective.  Thus
$X/S$ is smooth and proper.  It is projective: the tensor product of all
translates of an ample bundle, with its permutation linearization,
descends to an ample bundle.  In particular these examples are schemes,
not merely algebraic spaces.

\subsection{Both signs of the Hodge asymmetry}
\paragraph{D1 (Proved: both signs, and variation in a smooth family).}
At 57:287 Grothendieck asks for examples where the difference is
``$>0$, ou $<0$'' and where its terms vary with the parameter.
Let $k$ be algebraically closed of characteristic $2$, let $C,D$ be
elliptic curves, and let $u:D\to C$ be an isogeny.  On $A=C\times D$ put
\[
 g(x,y)=(-x+u(y),y),\qquad \Gamma=\langle g\rangle\simeq\mathbb Z/2.
\]
This is an involution because $g^2=1$.  Its tangent matrix and that of
its pullback action on $A^\vee=C^\vee\times D^\vee$ give
\begin{align*}
 \dim (t_A)_\Gamma&=2-\operatorname{rank}(du),\\
 \dim(t_{A^\vee})^\Gamma&=2-\operatorname{rank}(du^\vee),\\
 \dim(t_{A^\vee})^\Gamma-\dim(t_A)_\Gamma
   &=\operatorname{rank}(du)-\operatorname{rank}(du^\vee).
\end{align*}
Indeed the matrices of $g-1$ and $g^\vee-1$ in characteristic $2$ are
$\left(\begin{smallmatrix}0&du\\0&0\end{smallmatrix}\right)$ and
$\left(\begin{smallmatrix}0&0\\du^\vee&0\end{smallmatrix}\right)$.

Choose an ordinary elliptic curve $E/k$.  For
$u=F_{E/k}:E\to E^{(2)}$ its dual is Verschiebung $V$, with $dF=0$ and
$dV\ne0$.  The displayed difference is $-1$.  For $u=V$ the difference
is $+1$.  An explicit auxiliary elliptic curve is
\[
 Q:\quad y^2+xy=x^3+1,
 \qquad P=(0,1)\in Q[2](k)\setminus\{0\}.
\]
Take $B=Q$ and let the generator act by translation by $P$.  This is a
free action, so the construction above gives a smooth projective
threefold $X$.  Since $Y=Q/\langle P\rangle$ is elliptic,
\[
 h^1(X,\mathcal O_X)=1+\dim(t_{A^\vee})^\Gamma,
 \qquad
 h^0(X,\Omega_X^1)=1+\dim(t_A)_\Gamma.
\]
The first equality is the tangent space of \eqref{dgj:pic-tau}; the
second is étale descent of differential forms and the product formula
for $\Omega^1_{A\times Q}$.  It does not use the false assertion that
wild descent of $H^1(\mathcal O)$ is simply taking invariants.
The two examples have respectively $(h^{01},h^{10})=(2,3)$ and $(3,2)$.

For variation take $R=k[[t]]$ and the elliptic scheme
\begin{equation}\label{dgj:elliptic-family}
 E:\quad y^2+txy+y=x^3.
\end{equation}
Its discriminant is $1+t^3$, a unit.  The generic fibre is ordinary and
the special fibre is supersingular.  Repeating the construction with
$u=F$ or $V$ gives respectively $-1\to0$ and $+1\to0$ under
specialization.  The assertions about ordinarity may also be read
directly from the formal parameter $z=x/y$: multiplication by $2$ has
initial term $tz^2$, while at $t=0$ its first nonzero term has degree
$4$.  The symbolic certificate supplied with this paper checks the
discriminant and $[2](z)=tz^2\pmod {z^4}$.

\subsection{A nonflat fixed subgroup and nonflat $\operatorname{Pic}^{\tau}$}
\paragraph{D2 (Proved: the pathologies requested at 57:288).}
Keep \eqref{dgj:elliptic-family}, put $C=E^{(4)}$, $u=F^2:E\to C$, and
$A_0=C\times E$.  Let $g(x,y)=(-x+F^2y,y)$ and $h=-1$.  They commute,
are distinct involutions, and generate a faithful constant group
$\Gamma\simeq(\mathbb Z/2)^2$.  Scheme-theoretically
\[
 A_0^\Gamma=C[2]\times H,
 \qquad H=\ker(F^2:E\to C)\cap E[2].
\]
There is also a short independent check of the formal calculation.
For $z=x/y$ and $w=1/y$, the Weierstrass equation gives
$w+tzw+w^2=z^3$, hence $w\in z^3R[[z]]$.
Negation sends $(x,y)$ to $(x,y+tx+1)$, so
\[
 [-1](z)=\frac{z}{1+tz+w}=z+tz^2+O(z^3).
\]
Write the commutative formal group law as
$F(U,V)=U+V+cUV+d(U^2V+UV^2)+O((U,V)^4)$.
The identity $F(z,[-1](z))=0$ forces $c=t$.
Its diagonal has no cubic term in characteristic two, so
$[2](z)=tz^2\pmod {z^4}$.  This reasoning is independent of
rational-function duplication and the symbolic certificate.

The formal coordinate computation gives the exact local algebra
\[
 \mathcal O_H=R[z]/(z^4,tz^2).
\]
Here $\ker F^2$ is supported on zero with algebra $R[z]/(z^4)$, and
$[2](z)\equiv tz^2\pmod {z^4}$ supplies the additional equation.
Thus $H$ has generic rank $2$ and special rank $4$; equivalently the
nonzero classes $z^2,z^3$ are killed by $t$.  Since $C[2]$ has constant
rank $4$, $A_0^\Gamma$ has fibre lengths $8$ and $16$ and is not flat.

Take $A=A_0^\vee$ with the dual action, and $B=Q^2$, where the two
generators translate by $(P,0)$ and $(0,P)$.  The resulting $X/S$ is
smooth projective of relative dimension $4$, and
\[
 \operatorname{Pic}^{\tau}_{X/S}=Y^\vee\times A_0^\Gamma
\]
is not flat.  Nonflatness follows from the displayed local algebra:
if the product were flat over $S$, flatness would descend along its
smooth surjective projection to $A_0^\Gamma$.  This answers the existence question within the specified
quotient family itself.

\subsection{No abelian subscheme with the fibrewise Picard varieties}
\paragraph{D3 (Proved: failure of the stated specialization compatibility).}
At 57:288 the precise geometric request is the absence of a
``bon sous-schéma abélien'', followed by
``la formation de Pic variété n'est pas compatible avec la spécialisation''.
Use the same $A_0=C\times E$ and $g$, but only the cyclic group
$\langle g\rangle$, and take $A=A_0^\vee$, $B=Q$.  The fixed group
$G=A_0^g$ is the kernel of the $R$-homomorphism
\[
 C\times E\longrightarrow C,\qquad (x,y)\longmapsto[2]x-F^2y.
\]
On every geometric fibre its reduced neutral component is the image
of
\[
 n:E\longrightarrow C\times E,\qquad y\longmapsto(F^2y,[2]y).
\]
Indeed this is an elliptic image contained in $G$, and $G$ has
dimension one because its projection to $C$ is finite, a base change
of $F^2$.  But $\ker n=H$ has lengths $2$ and $4$ on the generic and
special fibres.

There cannot exist a closed abelian $R$-subscheme $P\subset G$ with
$P_s=(G_s^0)_{\mathrm{red}}$ for every geometric $s$.  If it existed,
$n$ would factor through $P$: this holds generically and extends
because $E$ is $R$-flat and $P$ is closed.  The resulting morphism
$E\to P$ would be a fibrewise isogeny of elliptic schemes and hence a
finite locally free morphism of constant degree.  Its kernel $H$
contradicts that constancy.

The same obstruction applies to
$\operatorname{Pic}^{\tau}_{X/S}=Y^\vee\times G$.  If a closed abelian
subscheme with the fibrewise reduced neutral components existed,
the morphism $Y^\vee\times E\to Y^\vee\times G$ induced by $(1,n)$
would factor through it and be a fibrewise isogeny of abelian schemes.
Its kernel again has nonconstant length.  This is an explicit, precise
failure of formation of the Picard variety to commute with this
specialization; it does not assert that its dimension changes.

\subsection{The product Igusa family and its limitations}
\paragraph{D4 (Product family and modular deformation analysis).}
For arbitrary elliptic schemes $E/S$ in characteristic $2$ and ordinary
$Q/S$ with a chosen nonzero étale $2$-torsion section, the usual Igusa
quotient $X=(E\times Q)/\langle(-1,+P)\rangle$ satisfies
\[
 \operatorname{Pic}^{\tau}_{X/S}=Y^\vee\times E^\vee[2].
\]
It is flat, and the fibrewise Picard variety is $Y^\vee$, compatibly
with base change.  Both first Hodge numbers are $2$.
The torsion Néron--Severi group is $\mathbb Z/2$ on an ordinary fibre
of $E$ and zero on a supersingular fibre, since it is the component
group of $E[2]$.  In the family \eqref{dgj:elliptic-family}, torsion
therefore disappears under specialization.  This product family is
not the full deformation family of the surface.  A first-order
nonflatness assertion is stated in accessible published literature,
but the proof audit below identifies a gap, and D4o subsequently refutes the assertion by computing the full ordinary deformation.  For the ordinary
Igusa surface $X_0$ put $I=H^1(X_0,T_{X_0})^\vee$ and
$R_1=k\oplus I$ with $I^2=0$, and take the first-order deformation
whose Kodaira--Spencer class is the identity of $H^1(X_0,T_{X_0})$.
Here $T_{X_0}=\mathcal O_{X_0}^{\oplus2}$ and
$h^1(\mathcal O_{X_0})=2$, so $\dim_k I=4$.
Chai--Oort, \S3.9(iv), printed p.~35
\cite{DGJ_ChaiOort2017}, states that
$\operatorname{Pic}^{\tau}_{X_1/R_1}$ is not flat.
The source identifies this assertion with Mumford's example.
The surface deformation
can be taken projective: the obstruction to lifting the square of
an ample line bundle is twice its square-zero obstruction, hence
zero in characteristic two, and ampleness over an Artinian base
is detected on the special fibre.

The annotated March 31, 1962 letter, p.~27, footnote~1
\cite{DGJ_GrothMumfordLetters}, gives the more detailed description
that the other connected component factors through a proper closed
subscheme of $\operatorname{Spec}R_1$.  Its auxiliary assertion
about \emph{every} nontrivial line bundle killed by two having
nonzero Hodge Chern class is too broad as printed: pull back a
nontrivial degree-zero two-torsion line bundle from the elliptic
Albanese $Y$.  The pullback remains nontrivial because the zero
section splits pullback, whereas its Hodge Chern class vanishes,
as degree zero line bundles on an elliptic curve have zero Chern
class in $H^1(\Omega^1)$.  We do not copy that overbroad auxiliary
claim into a new proof.  Direct Chai PDF downloads still return 403,
but the complete seven-page v2 has now been opened in web retrieval.
Theorem~2.4, pp.~6--7, and its Remark~(1), p.~7, have been inspected
in full.  The following caution applies to their final inference;
it does not by itself refute the actual Igusa statement; the independent proof D4o does so.

\paragraph{D4-gap (Proved caution on the inspected tangent argument).}
The proof argues that the off-neutral component cannot be a
flat torsor under the neutral component because their absolute
tangent dimensions are five and six, respectively, and their
tangent maps to the base have different ranks.  That inference
is false for torsors under a nonsmooth group.  Here is a complete
example with exactly those dimensions.  Over an algebraically
closed field $k$ of characteristic two put
\[
 R=k[e_1,e_2,e_3,e_4]/(e_i e_j:1\leq i,j\leq4),
 \qquad P=\operatorname{Spec}R[z]/(z^2-e_1).
\]
It is finite free of rank two over $R$, with basis $1,z$.
The unit $1+z$ satisfies $(1+z)^2=1+e_1$; multiplication on this
unit makes $P$ a $\mu_2$-torsor.  Explicitly, with $y^2=0$ the
coordinate of $\mu_2$, the torsor map is
$z_1\mapsto z$, $z_2\mapsto z+y+zy$.
On the two free rank-four coordinate modules over $R$ its
determinant is $1+e_1$, a unit.  This proves the torsor condition,
including faithful flatness.

The maximal ideal of $P$ has cotangent basis $z,e_2,e_3,e_4$,
because $e_1=z^2$.  Thus its tangent dimension is four, and its
tangent map to $\operatorname{Spec}R$ has image the hyperplane
$de_1=0$.  In contrast $\mu_2\times\operatorname{Spec}R$ has
tangent dimension five and a surjective tangent map to the base.
After multiplying both spaces by any smooth elliptic curve,
these dimensions become five and six, exactly as in the inspected
argument.  The first space remains a flat torsor under the
second group scheme.  Moreover $P(R)$ is empty, since every
principal unit of $R$ has square one, while $1+e_1\ne1$.
Consequently absence of an $R$-valued lift also does not imply
that $P$ factors through a proper closed subscheme of the base.
The finite algebra calculation is checked in
\texttt{check\_flat\_torsor.py} and its saved output.

The tangent inference is invalid. The complete ordinary calculation D4o below proves flatness for every Artinian deformation and therefore refutes the literal published first-order claim. This supersedes both the earlier Known classification and the subsequent unresolved verification status. The independent nonflat family D2 is unaffected.

\paragraph{A further source discrepancy and its independent replacement.}
Partsch's published classification \cite{DGJ_Partsch2013} gives, for the ordinary Igusa
case $p=d=2$, $h^1(T_X)=4$ in Proposition~6.1, p.~95,
whereas the Jacobian-deformation hull in Proposition~6.2, p.~99,
is $W[[t_E,q-1]][\sqrt q]$, whose relative tangent dimension is
two.  Theorem~4.4, p.~89, adds $h^1(C,\operatorname{Lie})=1$.
The claimed equality of the resulting tangent dimension with
four on p.~100 therefore needs correction or an additional
qualification.  We have not imported the classification to
close the calculation. The independent calculations D4o and D4s below
supply both miniversal hulls, including the quotient by changes of origin,
and prove flatness of the torsion Picard scheme for every Artinian
deformation. They were separately audited by root and fhi.
The explicit first-order deformation D4a also gives the direct
counterexample D4b to the universal lifting assertions on
pp.~102--106 of that article; the issue is therefore not only a
tangent-count discrepancy.

% BEGIN D4m independent modular degeneration
\paragraph{D4m (Proved: a modular Igusa degeneration with flat
$\operatorname{Pic}^{\tau}$ but incompatible Picard varieties).}
This example is independent of the full versal-ring calculation.
Let $k$ be algebraically closed of characteristic two, let $R=W(k)$,
and take both $E$ and $F$ to be the elliptic curve
\[
 Y^2Z+XYZ=X^3+XZ^2
\]
over $R$, with origin $O=[0:1:0]$ and $P=[0:0:1]$.
The discriminant is $-63$, a unit of $R$, so these are smooth elliptic
schemes. The point $P$ has order two and specializes to a nonzero
point; hence the special fibres are ordinary. Use subscripts to
distinguish $P_E$ and $P_F$.

Let $w\in\mathbb Z_2\subset R$ be the unique root congruent to zero
modulo two of
\[
 2w^2-w+8=0.
\]
It exists uniquely by Hensel's lemma, because the derivative is odd;
also $w\in8R$. The point
\[
 Q_F=[-2:1:w]\in F(R)
\]
lies on $F$, is fixed by elliptic negation
$[X:Y:Z]\mapsto[X:-Y-X:Z]$, and is nonzero on the generic fibre.
Thus it has order two generically and as an $R$-section, whereas
$(Q_F)_k=O$. It is the nontrivial generic section of the multiplicative
two-torsion subgroup. No choice of a root of a nonunit in $R$ is
being assumed here; the displayed Hensel equation defines $Q_F$.

Define two commuting involutions on $E\times_R F$ by
\[
 h(x,y)=(x+P_E,y+Q_F),\qquad
 \tau(x,y)=(-x,y+P_F).
\]
They generate the constant group $\Gamma=(\mathbb Z/2)^2$.
The action is free. For $h$ this follows from its nonzero $E$-translation;
for $\tau$ and $h\tau$ it follows from their $F$-translations
$P_F$ and $P_F+Q_F$, both of which specialize to a nonzero point.
The zero locus of either section is closed in $\operatorname{Spec}R$
and misses its closed point, so it is empty. Consequently
\[
 \mathcal X=(E\times_R F)/\Gamma
\]
is smooth projective over $R$, and its quotient map is finite etale.
For quotient existence one may use Stacks, Lemma~66.14.2, tag07S7;
an invariant ample tensor product descends to show projectivity.
Its special fibre is the ordinary Igusa surface
\[
 \mathcal X_k\simeq
 \bigl((E_k/\langle P_E\rangle)\times F_k\bigr)
 /\langle(-1,+P_F)\rangle.                              \tag{M1}
\]

\emph{The torsion Picard scheme is flat.}
First divide by $h$. Since $h$ is translation by the etale subgroup
$\langle(P_E,Q_F)\rangle$, the quotient
$A=(E\times F)/\langle h\rangle$ is an abelian scheme.
The copy of $F$ embeds in $A$, because that subgroup meets
$0\times F$ trivially. Its quotient is $\overline E=E/\langle P_E\rangle$.
Duality gives
\[
 0\longrightarrow\overline E^{\vee}\longrightarrow A^{\vee}
 \longrightarrow F^{\vee}\longrightarrow0.              \tag{M2}
\]
The linear part $\sigma$ of the remaining deck involution acts by
$-1$ on the first term and by $+1$ on the last term. If
$G=(A^{\vee})^{\sigma^{\vee}}$, the same elementary lifting argument
works over every base change: lift a point of $F^{\vee}$ to $A^{\vee}$,
then subtract half of its failure of invariance in
$\overline E^{\vee}$, taking the half fppf locally. Therefore
\[
 0\longrightarrow\overline E^{\vee}[2]\longrightarrow G
 \longrightarrow F^{\vee}\longrightarrow0               \tag{M3}
\]
is fppf exact. The map $G\to F^{\vee}$ is a finite locally free
torsor of degree four.

For the remaining double cover $A\to\mathcal X$, cyclic descent
gives
\[
 0\longrightarrow\mu_2\longrightarrow
 \operatorname{Pic}^{\tau}_{\mathcal X/R}
 \longrightarrow G\longrightarrow0.                    \tag{M4}
\]
Indeed an invariant algebraically trivial line bundle can be
linearized fppf locally: its chosen isomorphism with the deck
pullback has square a base unit, and adjoining a square root of
that unit corrects the cocycle. The ambiguity is exactly a
character of the cyclic group, namely $\mu_2$. The descended
class lies in $\operatorname{Pic}^{\tau}$ because its square is
the norm of an algebraically trivial bundle, hence is in the
neutral Picard component on each fibre. Conversely every torsion
Neron--Severi class pulls back into $A^{\vee}$, since the
Neron--Severi group of an abelian variety is torsion free.
These are fppf-sheaf statements on arbitrary $R$-schemes.
They also prove representability of the torsion Picard functor
here by finite-flat torsor descent from $G$.
Combining (M3) and (M4),
$\operatorname{Pic}^{\tau}_{\mathcal X/R}\to F^{\vee}$
is finite locally free of degree eight. In particular the full
torsion Picard scheme is flat over $R$.

\emph{The Picard varieties cannot specialize as an abelian subscheme.}
Let $f:F\to\mathcal X$ send $y$ to the class of $(0,y)$ and use
$f(0)$ as base point. For either geometric fibre $s$, every pointed
map $\mathcal X_s\to Z$ to an abelian variety pulls back to a
homomorphism $u_E+u_F:E_s\times F_s\to Z$. Invariance under $\tau$
forces $2u_E=0$ and $u_F(P_F)=0$. Since homomorphisms of abelian
varieties have no additive torsion, $u_E=0$. Invariance under $h$
then forces $u_F(Q_F)=0$. Conversely these conditions suffice for
descent. The Albanese variety and the map induced by $f_s$ are thus
\[
 \operatorname{Alb}(\mathcal X_s)=
 F_s/\langle P_F,Q_F\rangle.
\]
On the generic fibre the two distinct nonzero rational two-torsion
points $P_F,Q_F$ are independent, so the subgroup is all of $F[2]$
and the map has degree four. On the special fibre $Q_F=0$, so its
degree is two. Dualizing, restriction along $f$ on the fibrewise
Picard varieties has degrees
\[
 \deg f_\eta^*=4,\qquad \deg f_k^*=2.                    \tag{M5}
\]
If a closed abelian subscheme of
$\operatorname{Pic}^{\tau}_{\mathcal X/R}$ had these Picard varieties
as all its fibres, restriction along $f$ would be an isogeny of
abelian schemes with the two degrees in (M5). An isogeny of abelian
schemes is finite locally free, so its degree is locally constant.
This contradiction proves the asserted specialization failure
inside an actual modular Igusa degeneration, even though
$\operatorname{Pic}^{\tau}$ itself is flat.

The finite calculation records the discriminant, the projective
point and two-torsion identities, and Hensel approximations. See
\texttt{check\_igusa\_trait.py} and its saved output. The calculation
does not infer flatness from tangent dimensions and does not use
the full-hull assertion.

% END D4m independent modular degeneration

% BEGIN D4o completed ordinary calculation
% Completed ordinary calculation; independently audited by root and fhi.
\subsection*{D4o: the complete ordinary Igusa deformation calculation}
\textbf{Status: Proved.} The calculation was independently audited by root and fhi, including the opened Katz and Mumford--Fogarty sources. It covers ordinary first factors; the supersingular case is completed separately in D4s.

Let $k$ be algebraically closed of characteristic two, let $E_0,F_0$ be ordinary
elliptic curves, and let $P\in F_0[2](k)$ be nonzero. Put
\[
 A_0=E_0\times F_0,\quad \sigma_0=(-1,1),\quad
 \tau_0=t_{(0,P)}\sigma_0,\quad X_0=A_0/\langle\tau_0\rangle.
\]
The base category consists of local Artinian $W(k)$-algebras with residue
field $k$, and all deformations are marked on the closed fibre.

\paragraph{Step 1: lifting the cover and choosing its origin.}
The finite etale cover $A_0\to X_0$, together with its deck involution,
lifts uniquely to every deformation $X/R$, by invariance of the etale site
under a nilpotent immersion. This equivalence also applies to marked
isomorphisms: specify the identity lift on $A_0$ to remove the alternative
deck transformation. The lifted cover $A/R$ is smooth proper. A section
lifting $0\in A_0(k)$ exists by smoothness. It gives a unique abelian-scheme
structure, by Mumford--Fogarty, Proposition~6.15, pp.~124--126
\cite{DGJ_MumfordFogarty1982}, applied successively to small extensions.
This proposition assumes properness, not projectivity.

Write the deck map uniquely as $\tau=t_c\sigma$, where
$c=\tau(0)$ and $\sigma$ fixes the origin. Rigidity makes $\sigma$ a group
automorphism lifting $\sigma_0$, and $\sigma^2=1$. Thus the only additional
equation for the affine map is $(1+\sigma)c=0$.

Conversely every such triple $(A,\sigma,c)$ gives a free involution: its
fixed locus is closed in the proper scheme $A$ and has empty special
fibre, hence is empty. Its quotient is a smooth proper deformation of
$X_0$. At each Artinian stage $A$ is projective: a sufficiently high
power of an ample bundle on $A_0$ lifts through successive small
extensions, because their obstruction groups are killed by two.
An invariant ample bundle is obtained by tensoring a bundle with its
translate, so the quotient is projective too.

\paragraph{Step 2: Serre--Tate parameters and exact point coordinates.}
Choose bases $e_1,e_2$ of $T_2A_0(k)$ and $e_1^\vee,e_2^\vee$ of
$T_2A_0^\vee(k)$, with $e_2/2=P$. Use Katz's convention
\[
 q_{ij}=q(A/R;e_i,e_j^\vee)\in1+\mathfrak m_R.
\]
Katz, Theorem~2.1, pp.~148--149, identifies deformations with this matrix,
identifies duality with transposition, and gives the homomorphism-lifting
criterion. Therefore $\sigma_0$ lifts exactly when
\[
 q_{12}^2=q_{21}^2=1.                                      \tag{I1}
\]

Here is an $R$-point interpretation that avoids any unstated surjectivity
of a torsion subgroup. Choose a compatible system
$P_n\in F_0[2^{n+1}](k)$ with $2^nP_n=P$. For $n$ sufficiently large,
$[2^n]$ kills the formal group $\widehat A(\mathfrak m_R)$.
Choose an arbitrary lift $\widetilde P_n\in A(R)$ and put
\[
 c_0=[2^n]\widetilde P_n.
\]
It is independent of the lift and of sufficiently large $n$.
Indeed changes of lift lie in the killed formal group; replacing
$n$ by $n+1$ gives the same result because
$2\widetilde P_{n+1}-\widetilde P_n$ lies in that formal group.
It is functorial in $R$, reduces to $(0,P)$, and is fixed by $\sigma$,
because $\sigma_0P_n=P_n$. By Katz's explicit recovery of the
pushout homomorphism, pp.~150--152,
\[
 [2]c_0=(q_{21},q_{22})\in
 \widehat A(\mathfrak m_R)\simeq(1+\mathfrak m_R)^2.        \tag{I2}
\]
The isomorphism on the right is the one furnished by the chosen dual
Tate basis. It also identifies $\sigma$ on the formal group with
$(z_1,z_2)\mapsto(z_1^{-1},z_2)$.

Every lift $c$ of $(0,P)$ is uniquely $c_0+z$, with
$z=(z_1,z_2)\in(1+\mathfrak m_R)^2$. Consequently
\[
 (1+\sigma)c=(q_{21},q_{22}z_2^2).
\]
The affine involution condition is therefore exactly
\[
 q_{21}=1,\qquad q_{22}z_2^2=1.                           \tag{I3}
\]
Together with (I1), the functor of triples is represented by
\[
 R_{\rm tr}=W(k)[[t,s,u,v]]/(s^2+2s),                     \tag{I4}
\]
where
\[
 q_{11}=1+t,\quad q_{12}=1+s,\quad q_{21}=1,\quad
 q_{22}=(1+v)^{-2},\qquad z_1=1+u,\ z_2=1+v.
\]
This inverse in $q_{22}$ fixes Katz's convention; replacing $q$ by
$q^{-1}$ yields the earlier equivalent coordinate draft.

\paragraph{Step 3: the ring is a miniversal hull for the unpointed surface.}
Forgetting the chosen origin gives a formally smooth transformation
from the triple functor to $\operatorname{Def}_{X_0}$: over a small
extension, lift the cover uniquely and lift its prescribed origin by
smoothness. At first order this transformation is injective as well
as surjective. In fact any marked isomorphism of the unpointed covers
is a translation followed by a marked abelian isomorphism; rigidity
forces the latter to preserve the same $q$ parameters. Translation
by a formal point $d$ changes $c$ by $(1-\sigma)d$. In formal torus
coordinates this changes $z_1$ by a square and leaves $z_2$ fixed.
On $k[\epsilon]/(\epsilon^2)$ every principal unit has square one.
Thus all four tangent parameters in (I4) remain distinct. This proves
formal smoothness and tangent bijectivity, the defining hull property.
It also explains why the ring represents triples but need not represent
the surface deformation functor: higher-order changes of origin can
identify different triples.

The compatible formal family can be algebraized over
$\operatorname{Spec}R_{\rm tr}$. Let $\lambda_0$ be the product principal
polarization on $A_0$. The homomorphism $2\lambda_0$ lifts through all
stages, since its Serre--Tate condition is $q_{ij}^2=q_{ji}^2$, which
holds by (I1) and (I3). The canonical diagonal pullback of the
Poincare bundle attached to this symmetric homomorphism gives a
compatible relatively ample line bundle whose special-fibre class is
$4\lambda_0$. Formal existence algebraizes the abelian scheme and
the involution, its translation point, and the quotient. The hull is
unpolarized; the auxiliary polarization is used only for algebraization.

\paragraph{Step 4: the fixed dual group is flat.}
Let $\mathcal E^\vee,\mathcal F^\vee$ denote the elliptic deformations
with Serre--Tate coordinates $q_{11},q_{22}$. The condition $q_{21}=1$
and the transposition rule lift both the original inclusion and projection
to an exact sequence of abelian schemes
\[
 0\longrightarrow\mathcal E^\vee\longrightarrow A^\vee
 \stackrel{\rho}{\longrightarrow}\mathcal F^\vee
 \longrightarrow0.                                      \tag{I5}
\]
For precision, the inclusion condition is the $(1,2)$ entry of the
dual matrix, namely $q_{21}=1$; the projection has the same condition.
The compositions are zero by rigidity. The lifted inclusion is a closed
immersion and the projection is smooth with connected elliptic kernel,
as can be checked on the unique geometric fibre and by the smoothness
criterion; their kernel is therefore the stated elliptic scheme.
Rigidity shows that $\sigma^\vee$ acts by $-1$ on the kernel and by
$+1$ on the quotient.

Put $G=\ker(1-\sigma^\vee:A^\vee\to A^\vee)$.
The restricted map $G\to\mathcal F^\vee$ is an epimorphism of fppf
sheaves. To lift $b\in\mathcal F^\vee(T)$, first lift it fppf locally
to $a\in A^\vee$. The element $(1-\sigma^\vee)a$ lies in
$\mathcal E^\vee$. Divide it by two fppf locally in that elliptic
scheme and subtract the chosen half from $a$. The result is fixed.
Its kernel is exactly $\mathcal E^\vee[2]$. Hence
\[
 0\longrightarrow\mathcal E^\vee[2]\longrightarrow G
 \longrightarrow\mathcal F^\vee\longrightarrow0          \tag{I6}
\]
is fppf exact, and $G\to\mathcal F^\vee$ is a finite locally free
torsor of degree four. In particular $G/R$ is flat.

\paragraph{Step 5: the relative Picard scheme is flat.}
For the finite etale double cover $\pi:A\to X$, translation acts
trivially on $\operatorname{Pic}^0_{A/R}=A^\vee$. Pullback gives an
exact sequence of fppf sheaves
\[
 0\longrightarrow\mu_2\longrightarrow
 \operatorname{Pic}^{\tau}_{X/R}\stackrel{\pi^*}{\longrightarrow}
 G\longrightarrow0.                                    \tag{I7}
\]
Here are the descent details. An invariant relative class is represented
fppf locally by a line bundle $L$. Invariance of its relative class gives an isomorphism up to a base line bundle; trivialize that line bundle fppf locally to obtain
$\alpha:\tau^*L\simeq L$. The square of this isomorphism is a base
unit, since $\Gamma(A_T,\mathcal O)=\Gamma(T,\mathcal O)$ universally.
Adjoin a square root of that unit fppf locally and rescale $\alpha$;
it then obeys the cocycle condition, so $L$ descends to $X$.
The descended class $D$ satisfies
$D^{\otimes2}=\operatorname{Nm}_{A/X}(L)$ in the relative Picard
functor. Fibrewise the norm of a class in $A^\vee$ belongs to the
neutral Picard component, because the norm maps the connected
abelian variety $A^\vee$ to that component. Thus $D$ belongs to
$\operatorname{Pic}^{\tau}$.
Conversely a numerically trivial class pulls back into $A^\vee$, as
the Neron--Severi group of an abelian variety is torsion free.
The choices of linearization on the trivial bundle form precisely
$\operatorname{Hom}(\mathbb Z/2,\mathbb G_m)=\mu_2$, proving the
kernel assertion. These arguments work over arbitrary base changes
$T\to\operatorname{Spec}R$, not only at geometric points.

Thus $\operatorname{Pic}^{\tau}_{X/R}\to G$ is a finite locally free
$\mu_2$-torsor of degree two. Composed with (I6), it is finite locally
free of degree eight over $\mathcal F^\vee$, and is flat over $R$.
In particular the universal first-order ordinary Igusa deformation has
flat $\operatorname{Pic}^{\tau}$.
This refutes the literal conclusion of Chai's Theorem~2.4(iii),
while leaving its line-bundle obstruction and differing tangent-space
computations entirely compatible with flatness.

\paragraph{Opened source locators and independent review.}
Katz \cite{DGJ_Katz1981}, pp.~148--152, contains the coordinate,
duality, lift criterion, and pushout recovery used above. The finite
etale lifting statement is Stacks, Theorem~59.45.2, tag04DY,
\url{https://stacks.math.columbia.edu/tag/04DY}; quotient existence is
Lemma~66.14.2, tag07S7,
\url{https://stacks.math.columbia.edu/tag/07S7}. Both were opened.
Mumford--Fogarty, Proposition~6.15, pp.~124--126, was opened through
the primary-book OCR mirror listed in the bibliography.
The full seven-page Chai v2 was opened in web retrieval, and its final
flatness inference has no additional qualification beyond the tangent
comparison already tested by the explicit torsor example.

No classification beyond ordinary $E_0$ is claimed in this proof. Root and fhi independently checked (I2)--(I7), the origin quotient, and algebraization; fhi's audit is recorded separately. The flatness assertion for every Artinian ordinary deformation uses Steps1,2,4,5 alone and does not depend on miniversality.

\paragraph{Two concrete first-order checks.}
Over $k[\epsilon]/(\epsilon^2)$ the parameters $q_{12}=1+\epsilon$, all other $q_{ij}=1$, and $z_1=z_2=1$ give an actual nonproduct direction of the surface deformation, with flat torsion Picard scheme by (I6)--(I7). There is also a completely explicit product direction: take $E_\epsilon:y^2+xy=x^3+1+\epsilon$ and keep the ordinary translation factor fixed. Its nonidentity two-torsion component is $\operatorname{Spec}R[a]/(a^2-\epsilon)$, obtained by putting $x=0,y=1+a$. It is finite flat despite lacking an $R$-point. Its absolute tangent dimension is one smaller than that of the neutral $\mu_2$ component. Thus the nonsmooth-torsor behavior used to diagnose the published inference occurs in an Igusa family itself.

% END D4o completed ordinary calculation

% BEGIN D4s completed supersingular calculation
% Completed supersingular calculation; independently audited by root and fhi.
\subsection*{D4s: the supersingular Igusa hull and the canonical elliptic subgroup}
\textbf{Status: Proved.} Root and fhi independently audited the complete argument.
This section treats the remaining case in D4o. All bases in the
deformation calculation are local Artinian $W(k)$-algebras with residue
field an algebraically closed field $k$ of characteristic two; all
deformations and isomorphisms are marked on the closed fibre. Let
$E_0$ be supersingular, let $F_0$ be ordinary, and choose
$0\ne P\in F_0[2](k)$. Put
\[
 A_0=E_0\times F_0,\qquad \sigma_0=(-1,1),\qquad
 \tau_0=t_{(0,P)}\sigma_0,\qquad X_0=A_0/\langle\tau_0\rangle.
\]
The conclusions are
\[
 R_{X_0}^{\mathrm{hull}}
   =W(k)[[t,s,u,v]]/(s^4-ts^2-2s),                       \tag{S1}
\]
and flatness of $\operatorname{Pic}^{\tau}$ for every Artinian
deformation of $X_0$. The ring in (S1) is finite free of rank four over
$W(k)[[t,u,v]]$. Together with D4o, this gives the two cases of the full
Igusa calculation; the ordinary ring is
$W(k)[[t,s,u,v]]/(s^2+2s)$.

\paragraph{Step 1: the cover and the multiplicative subgroup.}
As in D4o Step1, every deformation $X/R$ uniquely lifts the marked
finite etale cover $A_0\to X_0$ and its deck involution. A chosen lift
of its origin makes $A/R$ an abelian scheme; write its deck involution
as $\tau=t_c\sigma$, where $\sigma$ is a group automorphism lifting
$\sigma_0$. Thus $\sigma^2=1$ and $(1+\sigma)c=0$.
The sources are Stacks tag04DY and Mumford--Fogarty,
Proposition6.15, pp.124--126, already opened in D4o.

Write $G=A[2^\infty]$ and $D=G^0$. The connected--etale sequence,
and Cartier duality applied to that sequence for $G^\vee$, give
\[
 G/D\simeq (\mathbb Q_2/\mathbb Z_2)_R,
 \qquad M:=((G^\vee)^{\mathrm{et}})^\vee
       \simeq\mu_{2^\infty}\ \subset D.                 \tag{S2}
\]
The marked special fibre identifies $M$ with the multiplicative
part of $F_0[2^\infty]$. These are canonical subgroups: the only
etale part of $G_0$ and of $G_0^\vee$ comes from $F_0$.
The quotient $G/M$ is a Barsotti--Tate group, being Cartier dual to
$(G^\vee)^0$. Its connected part is $D/M$, so $D/M$ is a connected
Barsotti--Tate group deforming $E_0[2^\infty]$.
For these facts one can use Tate, \S1.3--1.4, pp.160--161,
\S2.2, Proposition1 and equation(4), pp.162--164, and
\S2.3, p.166 \cite{DGJ_Tate1967}; the exact sequence and duality
are also obtained level by level from their finite flat versions.

The action of $\sigma$ on $M$ is $+1$, and its action on $D/M$ is
$-1$. Here no eigenspace splitting in characteristic two is being
assumed. A homomorphism between Barsotti--Tate groups over $R$ that
reduces to zero is zero: its image is killed by a fixed power of two
in the infinitesimal formal kernel, whereas its source is two-divisible.
This is the rigidity argument in Katz, Lemmas1.1.1--1.1.3,
pp.141--143, using the equivalence between connected Barsotti--Tate
groups and divisible formal groups. Applying it to $\sigma-1$ on
$M$ and $\sigma+1$ on $D/M$ proves the assertions.

\paragraph{Step 2: the canonical divisible system of actual points.}
Choose a compatible generator of $T_2F_0(k)$, or equivalently points
$e_r\in F_0[2^r](k)$ for $r\geq1$, with $e_1=P$ and
$2e_{r+1}=e_r$; put $e_0=0$. The formal group
$\widehat A(\mathfrak m_R)$ is killed by $2^N$ for sufficiently large
$N$ (Katz, Lemma1.1.1). For $r\geq0$ define
\[
 j(2^{-r})=[2^N]\widetilde e_{N+r}\in A(R),             \tag{S3}
\]
where $\widetilde e_{N+r}$ is any lift to $A(R)$, and take $N$ large.
Such lifts exist by smoothness. Changing the lift or increasing $N$
does not change (S3): the difference of two choices, or of
$2\widetilde e_{N+r+1}$ and $\widetilde e_{N+r}$, belongs to the
killed formal kernel. These identities also give
$2j(2^{-r-1})=j(2^{-r})$. Thus (S3) defines a functorial additive map
\[
 j:\mathbb Z[1/2]\longrightarrow G(R).
\]
Its values really are two-power torsion: their reductions are torsion
and a further uniform power of two kills the formal kernel. Every
value of $j$ is fixed by $\sigma$, since the chosen special-fibre
points are fixed and the same independence argument applies to their
lifts. Set $c_0=j(1/2)$ and $\kappa=j(1)=2c_0\in D(R)$.

Write $c=c_0+z$, where $z\in\widehat A(\mathfrak m_R)=D(R)$.
The last equality follows either from the formal-group description
of $D$ or directly by killing the formal kernel by a power of two.
The affine involution condition is
\[
 0=(1+\sigma)c=\kappa+(1+\sigma)z.
\]
Since $\sigma=-1$ on $D/M$, this forces
\[
                         \kappa\in M(R).                \tag{S4}
\]
This is the step that uses the existence of the affine deck map, not
merely a lift of its linear part.

\paragraph{Step 3: construction and uniqueness of the elliptic subgroup.}
As fppf sheaves over $R$ there is a presentation
\[
 G\simeq
 (D\oplus\underline{\mathbb Z[1/2]})/
       \{(-n\kappa,n):n\in\underline{\mathbb Z}\},       \tag{S5}
\]
whose map to $G$ is $(d,r)\mapsto d+j(r)$. This is a statement on
arbitrary test schemes $T/R$, not just on $R$-points. The image of
an element of $G(T)$ in the constant etale quotient is locally a
constant rational class in $\mathbb Z[1/2]/\mathbb Z$; choose a
representative $r$ locally and subtract the base section $j(r)_T$.
The difference lies in $D(T)$. The kernel consists exactly of the
pairs $(-n\kappa,n)$, which proves (S5) as a sheaf quotient.

By (S4) the analogous subquotient
\[
 H=(M\oplus\underline{\mathbb Z[1/2]})/
       \{(-n\kappa,n):n\in\underline{\mathbb Z}\}
       \ \subset G                                      \tag{S6}
\]
is defined. It fits in an exact sequence
$0\to M\to H\to\mathbb Q_2/\mathbb Z_2\to0$.
It is a Barsotti--Tate group, as can be checked without invoking an
extension-closure assertion: multiplication by two is surjective on
both end terms and therefore on $H$; for every $n$ the snake lemma
and $M/2^nM=0$ give
\[
 0\longrightarrow M[2^n]\longrightarrow H[2^n]
 \longrightarrow (\mathbb Z/2^n)_R\longrightarrow0.     \tag{S7}
\]
Thus $H[2^n]$ is a torsor under the finite flat group $M[2^n]$ over
the constant finite etale scheme $(\mathbb Z/2^n)_R$; it is finite
locally free of rank $2^{2n}$. These finite levels give a height-two,
dimension-one ordinary Barsotti--Tate group. Their monomorphisms into
$G[2^n]$ are closed immersions. Its special fibre is exactly the
given subgroup $F_0[2^\infty]\subset A_0[2^\infty]$.

Serre--Tate, in the general form of Katz Theorem1.2.1,
pp.143--144, now produces an elliptic scheme $F/R$ and a
homomorphism $F\to A$ lifting $F_0\hookrightarrow A_0$ and
realizing $H\hookrightarrow G$. This map is a closed immersion:
it is proper and quasi-finite, since its only geometric fibre is a
closed immersion, hence finite; the map on structure sheaves is
surjective by reduction and Nakayama. We identify $F$ with its image.
Both $M$ and $j$ are fixed by $\sigma$, so $\sigma|_F=1$ by
Serre--Tate full faithfulness. The quotient $E=A/F$ is an elliptic
scheme, and $\sigma$ induces $-1$ on it by rigidity.

This subgroup is unique among lifts of $F_0$ fixed pointwise by
$\sigma$. Indeed such a lift $F'$ has multiplicative subgroup $M$
and the same etale quotient. In (S3) one may choose all the lifts
inside the smooth scheme $F'$, so every $j(r)$ belongs to $F'$.
Consequently its Barsotti--Tate subgroup contains $H$; (S7) and
the ranks at every finite level force equality. Full faithfulness
then gives $F'=F$. In particular $F$, $E$, and the subsequent
classification are functorial in the deformation.

\paragraph{Step 4: an elementary graph description of $(A,\sigma)$.}
Let $\rho:A\to E=A/F$. The homomorphism $1-\sigma$ kills $F$, so
it factors uniquely as $a\rho$, where $a:E\to A$. Since
$\sigma=-1$ on $E$, one has $\rho a=[2]_E$.
The homomorphism
\[
 \Phi:E\times F\longrightarrow A,\qquad(e,f)\longmapsto a(e)+f
\]
is an fppf epimorphism with kernel
\[
 K=\{(e,h(e)):e\in E[2]\},\qquad
 h=-a|_{E[2]}:E[2]\longrightarrow F[2].                 \tag{S8}
\]
Indeed one can lift $\rho(x)$ by dividing it by two fppf locally in
$E$, and then the residual difference lies in $F$; the kernel formula
is immediate. Thus $\Phi$ is an isogeny of degree four.
On the special fibre $a(e)=(2e,0)$, so $h$ reduces to zero.
Its composite with $F[2]^{\mathrm{et}}\simeq\mathbb Z/2$ is
zero, by invariance of maps into an etale scheme under a nilpotent
immersion. Hence $h$ factors through $F[2]^0\simeq\mu_2$.
Cartier duality gives
\[
 \operatorname{Hom}(E[2],\mu_2)
  =\operatorname{Hom}(\mathbb Z/2,E^\vee[2])
  =E^\vee[2](R).                                       \tag{S9}
\]
Let $Q$ denote the corresponding point; it reduces to zero.

Conversely, elliptic deformations $E,F$ and such a point $Q$ define
$h_Q$ by (S9) and the abelian quotient
\[
 A_Q=(E\times F)/\operatorname{graph}(h_Q).               \tag{S10}
\]
The involution $(-1,+1)$ preserves the graph, because $-1=+1$ on
$E[2]$, and descends to $\sigma$. The factor $F$ embeds in $A_Q$,
and the quotient by it is $E$ via $[2]$ on the first factor.
The marked special-fibre identification is induced by
$(e,f)\mapsto(2e,f)$. The homomorphism $a$ is $e\mapsto[(e,0)]$,
so this construction recovers exactly (S8) and $Q$. By Step3 the
elliptic subgroup used in the reconstruction is unique. We have
therefore classified all the required $(A,\sigma)$, not just those
which already come with an elliptic fibration.

\paragraph{Step 5: the affine norm equation removes the $F$ parameter.}
The marked deformation space of $E_0$ has a parameter $t$.
Choose a coordinate $s$ at zero on the dual elliptic deformation,
and let $g_E(t,s)=0$ be its two-torsion equation. The ordinary
elliptic deformation $F$ has Serre--Tate parameter $q\in1+\mathfrak m$.
Thus the pairs $(A,\sigma)$ of Step4 have complete local base
\[
 B=W(k)[[t,q-1,s]]/(g_E(t,s)).                           \tag{S11}
\]
The marking identifies the formal group of $F$ with
$\widehat{\mathbb G}_m$. The canonical point $c_0$ in Step2 lies
in $F$, by its construction there, and Katz's pushout recovery,
pp.150--152, gives $[2]c_0=q$ in this formal-group coordinate.

Choose a formal-variety section $b(u)$ of
$\widehat A\to\widehat E$ over $B$, through the identity; this
exists by formal smoothness and a choice of relative formal
coordinates. It is not required to be a homomorphism. Every formal
point of $A$ can then be written uniquely as $b(u)+(1+v)$, with
$1+v$ in the formal group of $F$. Since $\sigma=-1$ on $E$,
write $(1+\sigma)b(u)=n(t,q,s,u)$ in the formal group of $F$.
It is a principal unit and satisfies $n(t,q,s,0)=1$ identically.
The affine condition becomes exactly
\[
                 q\,n(t,q,s,u)(1+v)^2=1.               \tag{S12}
\]
The derivative of its left side minus one with respect to $q$ is
one at the closed origin. The formal implicit-function theorem
therefore eliminates $q$ uniquely. This elimination is valid over
the possibly nonreduced base (S11): lift $n-1$, which is divisible
by $u$, to a power series before imposing $g_E$, eliminate $q$
there, and then impose $g_E(t,s)$, which is independent of $q$.
The functor of triples $(A,\sigma,c)$ is consequently represented
by
\[
                    W(k)[[t,s,u,v]]/(g_E(t,s)).          \tag{S13}
\]

Forgetting the chosen origin is formally smooth onto
$\operatorname{Def}_{X_0}$, by the cover-lifting and smooth-origin
argument in D4o. It is tangent bijective: a marked isomorphism of
the covers is a translation by a formal point $d$, followed by a
marked group isomorphism. The group isomorphism preserves the
classified data; the translation changes $c$ by $(1-\sigma)d$.
On $k[\epsilon]/(\epsilon^2)$ this last expression is zero,
because the derivative of $\sigma_0=(-1,1)$ is the identity in
characteristic two. Hence no tangent parameters in (S13) are
identified. This proves the hull property. As before, the assertion
is miniversality for unpointed surfaces, not representability of that
functor; changes of origin can act at higher orders.

\paragraph{Step 6: the exact supersingular equation.}
Choose $E_0:y^2+y=x^3$, which represents the unique supersingular
elliptic isomorphism class over $k$. A universal marked deformation
is
\[
                  E_t:\ y^2+txy+y=x^3,
                  \qquad\Delta=t^3-27.                 \tag{S14}
\]
For completeness this normal form follows directly from general
Weierstrass equations. At the displayed special equation, changes
$x=\lambda^2x'+r$, $y=\lambda^3y'+\lambda^2 b x'+d$ have an
invertible Jacobian on the four coefficients
$(a_2,a_3-1,a_4,a_6)$: modulo two its diagonal entries, in the
order $(r,\lambda-1,b,d)$, are $(3,-3,-1,-1)$.
Thus these four coefficients can be normalized formally to
$(0,0,0,0)$, leaving $a_1=t$. There are no infinitesimal pointed
automorphisms of an elliptic curve. This proves the asserted
one-parameter marked deformation chart. The canonical principal
polarization identifies $E_t^\vee$ with $E_t$.

Near $O=[0:1:0]$, use $s=X/Y=x/y$ and $w=Z/Y=1/y$.
The curve equation and inversion are
\[
 w+t s w+w^2=s^3,\qquad
 (s,w)\longmapsto
 \left(-\frac{s}{1+ts+w},-\frac{w}{1+ts+w}\right).       \tag{S15}
\]
There is a unique formal solution $w=w_t(s)\in(s^3)$.
Because the denominator in (S15) is a unit and $w\in(s)$, the
fixed-point ideal of inversion is generated by
$s(2+ts+w)$. The following exact identity in the curve ring removes
the auxiliary coordinate:
\[
 (w-1)(w+2+ts)=s^3-ts-2.
\]
Since $w-1$ is a unit at the origin, the two-torsion ideal is
\[
 (s(2+ts+w))=(s^4-ts^2-2s).                             \tag{S16}
\]
Equations (S13) and (S16) prove (S1). All of the supersingular
two-torsion lies in this formal neighborhood. The polynomial is
monic in $s$ and has its lower coefficients in the maximal ideal of
$W(k)[[t,u,v]]$; Weierstrass division makes the complete quotient
free with basis $1,s,s^2,s^3$. At $2=t=0$ its two-torsion factor is
$k[s]/(s^4)$, displaying the fourth-order thickening explicitly.

\paragraph{Step 7: algebraization and flatness of the torsion Picard scheme.}
The hull has an algebraic realization over its complete local base.
The elliptic deformation $E_t$ is given by (S14); the ordinary
elliptic family $F_q$ is effective using its Weierstrass deformation
chart and the Serre--Tate coordinate isomorphism. The point $Q$ is
an actual section of its finite two-torsion scheme. Hence (S10) is
an abelian scheme over the complete ring after substituting the
solution of (S12). The compatible point $c$ is effective as well:
choose an affine neighborhood of its closed-fibre point in $A$;
the compatible Artinian sections give its coordinate values in the
complete base ring. The resulting affine involution is free, because
its proper fixed locus has empty closed fibre.

There is a useful explicit ample bundle, so no unpolarized formal
existence assumption is needed. Since the kernel of
$\Phi:E\times F\to A$ lies in $(E\times F)[2]$, multiplication
by two factors as $g\Phi$ for an isogeny $g:A\to E\times F$.
Pull back the product ample bundle by $g$, then tensor it with its
$\tau$-pullback. Swapping the factors supplies a $\tau$-linearization.
It descends to an ample bundle on the finite etale quotient $X$.
This gives the claimed algebraized projective hull family.

Finally the canonical subgroup constructed in Step3 gives
\[
 0\longrightarrow E^\vee\longrightarrow A^\vee
 \longrightarrow F^\vee\longrightarrow0,
\]
with $\sigma^\vee=-1$ on the kernel and $+1$ on the quotient.
Put $G_\sigma=\ker(1-\sigma^\vee)$. Exactly as in D4o Steps4--5,
fppf division by two in $E^\vee$ and cyclic line-bundle descent give
\[
 0\longrightarrow E^\vee[2]\longrightarrow G_\sigma
 \longrightarrow F^\vee\longrightarrow0,
 \qquad
 0\longrightarrow\mu_2\longrightarrow
 \operatorname{Pic}^{\tau}_{X/R}\longrightarrow G_\sigma
 \longrightarrow0.                                    \tag{S17}
\]
Both surjections are torsors in the fppf topology. Their degrees
are four and two. In the second sequence an invariant relative
class is represented fppf locally by a line bundle; trivialize the
base factor in its invariance isomorphism, then adjoin a square root
of the cocycle unit to linearize it. The descended class is torsion
modulo the neutral component by the norm identity, just as in (I7).
These steps do not require any section of a nonsmooth torsor over
the original base. Thus $\operatorname{Pic}^{\tau}_{X/R}$ is finite
locally free of degree eight over $F^\vee$ and is flat over $R$.
This conclusion holds for every Artinian deformation already by
Steps1--3 and (S17), independently of the explicit hull calculation.

\paragraph{Finite certificate and scope.}
The code \texttt{check\_igusa\_supersingular.py} records the exact
Weierstrass discriminant, the inversion identities, the polynomial
identity in (S16), the normal-form Jacobian, and the special-fibre
length four. Its output does not purport to mechanize the
Barsotti--Tate construction, fppf descent, or the hull argument.
Those are supplied in full above. Root and fhi independently checked the subgroup construction, the
classification, the norm elimination, the hull and effectivity arguments,
and the two torsor sequences. Their audits are recorded in the
check-event ledger.

% END D4s completed supersingular calculation

\subsection{No compatible Picard variety over dual numbers}
\paragraph{D4a (Explicit counterexample to Artinian compatibility).}
Let $k=\overline{\mathbb F}_2$, $R=k[\epsilon]/(\epsilon^2)$,
and take two labelled copies $E_0,F_0$ of the ordinary elliptic
curve
\[
 y^2+xy=x^3+1,
 \qquad P=(0,1)\in F_0[2](k).
\]
Its discriminant is one and the coefficient of $xy$ is one, so it
is smooth ordinary; inversion sends $(x,y)$ to $(x,y+x)$,
showing that the displayed nonzero point has order two.
Set $A_0=E_0\times F_0$ and
$X_0=A_0/\langle(-1,t_P)\rangle$, as in D4o.
Choose the marked abelian deformation $A/R$ with Katz coordinates
\[
 (q_{ij})=\begin{pmatrix}1&1+\epsilon\\1&1\end{pmatrix}.
                                                        \tag{D4a.1}
\]
The involution and its translation lift: in D4o, equations
(I1)--(I3), choose $z_1=z_2=1$. Indeed
$(1+\epsilon)^2=1$ and the second row is $(1,1)$.
This gives an actual smooth projective quotient deformation
$X/R$ and a finite étale double cover $\pi:A\to X$.
Its torsion Picard scheme is flat, by D4o.

Nevertheless there is \emph{no} abelian $R$-scheme $B$ with a
homomorphism
\[
 \iota:B\longrightarrow\operatorname{Pic}^{\tau}_{X/R}
                                                        \tag{D4a.2}
\]
whose special fibre identifies $B_k$ isomorphically with
$\operatorname{Pic}^{\rm var}_{X_0/k}$ and whose special-fibre
map is its natural inclusion. In particular no abelian closed
subscheme of the relative Picard scheme can have the specified
reduced identity component on the closed fibre. The assertion
rules out existence, not merely a preferred construction.

Here is the proof. Put $Y_0=F_0/\langle P\rangle$ and write
$\phi:F_0\to Y_0$ for the étale degree-two isogeny. The product
Igusa calculation gives
\[
 \operatorname{Pic}^{\tau}_{X_0/k}
     \simeq Y_0^\vee\times E_0^\vee[2],
 \qquad \operatorname{Pic}^{\rm var}_{X_0/k}=Y_0^\vee.
\]
Under pullback to $A_0$, the latter inclusion becomes
\[
 f_0=(0,\phi^\vee):Y_0^\vee
                         \longrightarrow E_0^\vee\times F_0^\vee.
                                                        \tag{D4a.3}
\]
This also follows directly by pulling back a degree-zero line
bundle from the quotient map $X_0\to Y_0$.
If (D4a.2) existed, compose it with $\pi^*$ to the relative
Picard scheme of $A$. This composite lands in its open neutral
component $A^\vee$: it does so on the closed fibre, and an
Artinian nilpotent thickening has the same underlying topological
space, so the inverse image of the complementary open is empty.
We therefore obtain a homomorphism of ordinary abelian $R$-schemes
\[
 f:B\longrightarrow A^\vee
\]
lifting $f_0$. Both source and target are abelian schemes, so its
dual $f^\vee:A\to B^\vee$ is a homomorphism as well; no image,
kernel-flatness, or splitting assertion is needed.

Choose the ordinary étale Tate bases for $A_0$ and $A_0^\vee$
as in D4o, with indices $E,F$. Since $\ker\phi^\vee=\mu_2$,
its map on étale $2$-adic Tate modules is an isomorphism. Choose
the generator $x\in T_2Y_0^\vee(k)$ so that
\[
 T_2(f_0)(x)=e_F^\vee.
\]
The dual of (D4a.3) is $(0,\phi):E_0\times F_0\to Y_0$,
so
\[
 T_2(f_0^\vee)(e_E)=0.
\]
Katz's necessary and sufficient morphism-lifting criterion
\cite[Theorem 2.1(4), p.~149]{DGJ_Katz1981}, applied to
$f:B\to A^\vee$, would imply
\[
 q_B\bigl(x,T_2(f_0^\vee)e_E\bigr)
    =q_{A^\vee}\bigl(T_2(f_0)x,e_E\bigr).
\]
The left side equals one. By the transposition formula for dual
abelian schemes, Katz's Theorem~2.1(3), pp.~148--149, the right
side equals
\[
 q_{A^\vee}(e_F^\vee,e_E)=q_A(e_E,e_F^\vee)
      =q_{12}=1+\epsilon\ne1.
\]
This contradiction proves the assertion. It uses the exact
Serre--Tate lifting criterion on the special-fibre Tate modules,
not a claim that an arbitrary point lifts through a torsion
isogeny. It remains valid for every possible elliptic deformation
$B$ of $Y_0^\vee$: its first pairing entry is evaluated at zero
in the second slot, independently of its own parameter.

The leaf 57:311 asks ``Existence de'' $\operatorname{Pic}^{00}$
(the superscript is transcribed uncertainly as $00$ or $\infty$),
``dans le cas Artinien'' with ``local'' uncertain. Under the
compatible-abelian-subscheme interpretation specified above,
the answer is negative already on an Artinian base of length two. The full torsion Picard scheme
of this same example is flat. The two properties must not be
identified. The calculation of the matrix and the nonzero
obstruction in $R$ is certified in
\texttt{check\_artin\_picvar.py}; existence and the contradiction
are the geometric proof above.

The targeted independent audits by root and bce checked the étale Tate
isomorphism for $\phi^\vee$, the zero dual entry, the exact
transposition convention, and the existence of the quotient from
D4o. No gap was found. This explicit Serre--Tate proof is consistent
with Grothendieck's published first-order Igusa remark
\cite[Remark 3.4, p.~238]{DGJ_FGAVI1962}, not a claim of priority.

% BEGIN D4b
\paragraph{D4b (Refuted universal lifting of the Albanese fibration).}
\emph{Existing-source check.}
Grothendieck's earlier first-order Igusa obstruction to a compatible
abelian Picard subgroup is already recorded in
\cite[Remark 3.4, p.~238]{DGJ_FGAVI1962}.
The exact lifting assertions addressed here were read in the published
2013 text of \cite{DGJ_Partsch2013}, pp.~102--106.
No correction to these assertions was found in the retrieved
publisher and author-source search; this limited search is not a
claim that no correction exists. No priority claim is made for the
obstruction. Its fully specified deformation is D4a.

The explicit deformation $X/R$ of D4a, with
$R=\overline{\mathbb F}_2[\epsilon]/(\epsilon^2)$, does not
admit a lifting of the Albanese elliptic fibration
\[
 f_0:X_0\longrightarrow Y_0=F_0/\langle P\rangle.
\]
Here a lifting means a morphism $f:X\to C$ over $R$, where
$C/R$ is a smooth proper deformation of $Y_0$, and with
special fibre identified with $f_0$. The special fibration is
Jacobian: its section is induced by $\{0\}\times F_0$ in the
cover $E_0\times F_0$.

Suppose such $f$ exists. The relative Picard scheme
$B=\operatorname{Pic}^0_{C/R}$ is an elliptic abelian scheme,
and pullback is a homomorphism of $R$-group schemes
\[
 f^*:B\longrightarrow\operatorname{Pic}_{X/R}.
\]
Its image lies in $\operatorname{Pic}^{\tau}_{X/R}$, since
pullback preserves algebraic triviality on geometric fibres.
Its special fibre is precisely the canonical inclusion
$Y_0^\vee=\operatorname{Pic}^{\rm var}(X_0)
\hookrightarrow\operatorname{Pic}^{\tau}(X_0)$, because $f_0$
is the Albanese map. This is the homomorphism ruled out by
D4a, for every possible deformation $B$. Thus $f$ cannot
exist. No assertion about the scheme-theoretic image of
$f^*$, or its flatness, enters this contradiction.

This refutes the literal universal lifting assertions in
\cite[Proposition 6.6, pp.~102--103; Theorem 6.7, p.~104]{DGJ_Partsch2013}:
the opened published text states that every deformation of
a Jacobian bielliptic surface extends its elliptic fibration,
and then claims the same for the Albanese fibration of every
bielliptic surface. The special fibre above is exactly a
Jacobian ordinary Igusa surface, and D4a constructs an actual
smooth projective deformation. The stronger assertion that
both fibrations lift, in Proposition~6.9, p.~106, is refuted
by the same example. This does not make a claim about other
cases of that article's classification.

The last trace step of the proof on p.~103 illustrates a
specific danger. Killing a connected subgroup and inducing
zero on an etale quotient does not force the original map to
vanish: its remaining map may go from that etale quotient to
a connected infinitesimal target. Over this very $R$ there
is a nontrivial homomorphism
\[
 h:(\mathbb Z/2)_R\longrightarrow\mu_{2,R},\qquad
 h(0)=1,\quad h(1)=1+\epsilon,
\]
whose special fibre is zero. The relation
$(1+\epsilon)^2=1$ verifies the group law. Thus rigidity of
maps between etale group schemes cannot be applied when
only the source is etale. The complete geometric refutation
above uses D4a; this finite-group example identifies an
invalid inference without being substituted for that proof.
The Hopf-map and group-law identities are recorded by
\texttt{check\_partsch\_etale\_step.py} and its saved output.

% END D4b
